% =============================================================================
% Semana 06 — Trabalho 3 + Correlação: Pearson, Scatter e Causalidade
% Aula de 3 horas (19h–22h) | 18/03/2026
% =============================================================================

\chapter{Semana 06 (18/03/2026) — Correlação, Scatter e Causalidade}

% ---------------------------------------------------------------------------
\section{Objetivo da semana}
% ---------------------------------------------------------------------------

Pessoal, até aqui analisamos \emph{uma variável de cada vez} (análise
univariada). Hoje damos um passo fundamental: vamos olhar para \textbf{dois
atributos numéricos ao mesmo tempo} e perguntar — eles se movem juntos?
Quando as vendas sobem, o tempo de entrega também sobe? À medida que o
salário aumenta, a satisfação aumenta?

Essa é a análise bivariada, e a ferramenta central é o \textbf{coeficiente
de correlação de Pearson}. E vou adiantar a lição mais importante da aula:
\textbf{correlação não é causalidade} — e confundir os dois já causou
bilhões de reais em decisões erradas.

\begin{BoardBox}
\textbf{Roteiro da aula (19h–22h)}
\begin{enumerate}
  \item 19h–20h — \textbf{Apresentações:} Trabalho 3 (diagnóstico de dataset)
  \item 20h–20h30 — \textbf{TEORIA:} Scatter plot — ler padrões, clusters e outliers visuais
  \item 20h30–21h — \textbf{TEORIA:} Correlação de Pearson — fórmula, escala e interpretação
  \item 21h–21h30 — \textbf{TEORIA:} Significância estatística — amostra mínima e p-valor
  \item 21h30–21h40 — \textbf{TEORIA:} Correlação $\neq$ Causalidade + exemplos clássicos
  \item 21h40–21h55 — \textbf{PRÁTICA EM DUPLA:} Tabelas 1 e 2
  \item 21h55–22h — \textbf{ATIVIDADE INDIVIDUAL FINAL}
\end{enumerate}
\end{BoardBox}

\begin{NoteBox}
\textbf{Leitura base:} Cap.\ 03 (Correlação) + Introdução ao Cap.\ 04 (scatter).
Livro: \textit{Estatística Prática para Cientistas de Dados} (Bruce).
\end{NoteBox}

% =============================================================================
\section{Parte I — Apresentações: Trabalho 3 (19h–20h)}
% =============================================================================

\begin{BoardBox}
\textbf{Checklist de avaliação T3 — Diagnóstico de dataset:}
\begin{tabular}{|p{4.5cm}|p{8.5cm}|}
\hline
\textbf{Ponto} & \textbf{O que observar} \\
\hline
Dataset adequado & $\geq 50$ linhas, $\geq 3$ colunas numéricas \\
Medidas de posição & Calculadas para todas as colunas numéricas \\
Medidas de dispersão & DP + IQR + identificação de outliers por Tukey \\
Missing tratado & Estratégia documentada (remover, imputar, flag) \\
Vizualizações & Boxplot + histograma com interpretação \\
Big Idea & 1 frase com dado, significado e decisão \\
\hline
\end{tabular}
\end{BoardBox}

% =============================================================================
\section{Parte II — Teoria: Correlação (20h–21h40)}
% =============================================================================

% ---------------------------------------------------------------------------
\subsection{20h–20h30: Scatter Plot — lendo o gráfico de dispersão}
% ---------------------------------------------------------------------------

\begin{BoardBox}
\textbf{O que é o scatter plot (para o quadro)}

O scatter plot (gráfico de dispersão) coloca dois atributos nos eixos —
X e Y — e plota cada observação como um ponto. O padrão visual dos pontos
revela a relação entre as variáveis.

\textbf{O que observar em um scatter plot:}
\begin{itemize}
  \item \textbf{Direção:} os pontos sobem da esquerda para a direita?
  (relação positiva) Ou descem? (relação negativa) Não há padrão? (sem relação)
  \item \textbf{Intensidade:} os pontos estão concentrados ao longo de uma
  linha imaginária (forte) ou espalhados aleatoriamente (fraca)?
  \item \textbf{Forma:} é linear? Curvilínea (ex.: crescimento exponencial)?
  \item \textbf{Clusters:} há grupos distintos de pontos? (indicativo de subgrupos)
  \item \textbf{Outliers:} pontos muito afastados do padrão geral — investigar!
  \item \textbf{Cor por categoria:} pintar os pontos por uma terceira variável
  (ex.: região, tipo de cliente) pode revelar que a relação só existe em um subgrupo.
\end{itemize}
\end{BoardBox}

\begin{SolvedBox}
\textbf{Exemplos de padrões em scatter plots:}
\begin{itemize}
  \item \textbf{Relação positiva forte:} horas de estudo × nota na prova.
  \item \textbf{Relação positiva fraca:} temperatura × vendas de sorvete
  (real, mas com muita variação).
  \item \textbf{Relação negativa forte:} preço do produto × volume de vendas.
  \item \textbf{Sem relação:} número do sapato × salário.
  \item \textbf{Curvilínea:} dose de medicamento × resposta (sobe até um ponto,
  depois cai — Pearson não captura isso!).
  \item \textbf{Clusters:} scatter de IMC × pressão arterial separado por fumantes
  e não-fumantes — revela dois padrões distintos no mesmo gráfico.
\end{itemize}
\end{SolvedBox}

% ---------------------------------------------------------------------------
\subsection{20h30–21h: Correlação de Pearson}
% ---------------------------------------------------------------------------

\begin{FormulaBox}{Coeficiente de correlação de Pearson ($r$)}
\[
r = \frac{\sum_{i=1}^{n}(x_i - \bar{x})(y_i - \bar{y})}
         {\sqrt{\sum_{i=1}^{n}(x_i - \bar{x})^2} \cdot \sqrt{\sum_{i=1}^{n}(y_i - \bar{y})^2}}
\]
\begin{itemize}
  \item $r = +1$: correlação positiva perfeita (toda vez que X sobe, Y sobe na mesma proporção)
  \item $r = -1$: correlação negativa perfeita
  \item $r = 0$: sem correlação linear (pode ainda haver relação não-linear!)
  \item Interpretação prática: $|r| < 0{,}3$ fraca; $0{,}3 \leq |r| < 0{,}7$ moderada;
  $|r| \geq 0{,}7$ forte
\end{itemize}
\end{FormulaBox}

\begin{SolvedBox}
\textbf{Cálculo passo a passo de Pearson com 5 pontos:}

Dados: horas de treino por semana (X) × performance (Y, escala 0–100)

\begin{center}
\begin{tabular}{|c|r|r|r|r|r|}
\hline
Aluno & $x_i$ & $y_i$ & $(x_i-\bar{x})$ & $(y_i-\bar{y})$ & $(x_i-\bar{x})(y_i-\bar{y})$ \\
\hline
A & 2 & 50 & $-4$ & $-20$ & $80$ \\
B & 4 & 60 & $-2$ & $-10$ & $20$ \\
C & 6 & 70 & $0$ & $0$ & $0$ \\
D & 8 & 80 & $+2$ & $+10$ & $20$ \\
E & 10 & 90 & $+4$ & $+20$ & $80$ \\
\hline
\textbf{Média} & $\bar{x}=6$ & $\bar{y}=70$ & & & \textbf{Soma} $= 200$ \\
\hline
\end{tabular}
\end{center}

$\sum(x_i - \bar{x})^2 = 16+4+0+4+16 = 40$\\
$\sum(y_i - \bar{y})^2 = 400+100+0+100+400 = 1000$

\[
r = \frac{200}{\sqrt{40} \cdot \sqrt{1000}} = \frac{200}{\sqrt{40.000}} = \frac{200}{200} = 1{,}0
\]

Correlação perfeita — porque os dados foram criados exatamente assim.
Na prática, $r = 0{,}85$ já é muito forte e raro.
\end{SolvedBox}

\begin{BoardBox}
\textbf{Limitações do coeficiente de Pearson (para o quadro)}
\begin{itemize}
  \item Mede apenas \textbf{relação linear}. Uma relação curvilínea forte pode ter
  $r \approx 0$.
  \item É fortemente influenciado por \textbf{outliers}. Um único ponto extremo
  pode elevar ou destruir a correlação calculada.
  \item Não funciona bem com \textbf{variáveis ordinais} — use Spearman nesse caso.
  \item Variáveis com distribuições fortemente assimétricas: transformação logarítmica
  antes de calcular Pearson.
  \item \textbf{Sempre olhe o scatter plot} antes de confiar no número de Pearson —
  o Quarteto de Anscombe (1973) mostra quatro datasets com $r$ idêntico e padrões
  completamente diferentes.
\end{itemize}
\end{BoardBox}

% ---------------------------------------------------------------------------
\subsection{21h–21h30: Significância Estatística — Amostra Mínima e p-valor}
% ---------------------------------------------------------------------------

Até agora vimos como calcular e interpretar o coeficiente de Pearson.
Mas será que todo $r$ calculado merece confiança?
Nesta parte, vamos entender por que o tamanho da amostra importa tanto
quanto o valor de~$r$, e como usar o \textbf{p-valor} para decidir se
uma correlação é real ou mero acaso.
A ideia central é simples: um número sozinho não basta — precisamos
saber se ele se sustenta quando aumentamos ou mudamos os dados.

\begin{BoardBox}
\textbf{O problema: $r$ sozinho pode ser enganoso (para o quadro)}

Imaginem dois cenários:
\begin{itemize}
  \item \textbf{Cenário A:} analisei 500 clientes e achei $r = 0{,}35$.
  \item \textbf{Cenário B:} analisei 5 clientes e achei $r = 0{,}88$.
\end{itemize}

Qual achado é mais confiável? O segundo parece mais impressionante,
mas baseia-se em \textbf{apenas 5 observações} — qualquer coincidência
aleatória pode gerar um $r$ alto. O primeiro é modesto, mas com 500
observações é quase impossível que seja acaso.

Para responder ``esse $r$ é real ou pode ser pura coincidência?'',
precisamos de um \textbf{teste de significância} — e é aí que entra
o \textbf{p-valor}.
\end{BoardBox}

Percebam que o problema central é a \emph{quantidade} de dados.
Com poucos pontos, qualquer padrão pode ser coincidência.
Isso nos leva à próxima pergunta natural: afinal, quantos dados
são suficientes para confiar no resultado?

\begin{SolvedBox}
\textbf{A dúvida natural: ``quantos dados eu preciso?''}

Ao ver o cenário de 500 vs 5 observações, a primeira pergunta que
surge é: \textit{``Existe um número mínimo para eu confiar no meu resultado?''}

A resposta depende de três coisas:
\begin{enumerate}
  \item \textbf{Nível de confiança:} quanta certeza você quer? (geralmente 95\%)
  \item \textbf{Margem de erro:} quanta imprecisão você tolera?
  \item \textbf{Variabilidade dos dados:} quanto os dados oscilam?
\end{enumerate}

Antes de chegar à fórmula, precisamos entender o que significa
esse ``95\% de confiança'' — porque é o mesmo conceito que vai aparecer
no teste de significância logo a seguir.
\end{SolvedBox}

Para responder ``quantos dados eu preciso?'', primeiro temos que
entender um conceito fundamental: o \textbf{nível de confiança}.
Ele aparece em toda análise estatística e vai ser a ponte entre
amostragem e significância.
Quando alguém diz ``pesquisa com 95\% de confiança'', o que exatamente
isso quer dizer?

\begin{BoardBox}
\textbf{Nível de confiança — o que significa 95\%? (para o quadro)}

Imaginem que eu repita a mesma pesquisa 100 vezes, com 100 amostras
diferentes da mesma população. Se eu usar nível de confiança de 95\%:

\begin{itemize}
  \item Em \textbf{95 dessas 100 pesquisas}, o intervalo que eu calcular
  vai conter o valor real da população.
  \item Em \textbf{5 dessas 100 pesquisas}, por azar, o intervalo vai
  errar — não vai conter o valor real.
\end{itemize}

\textbf{Tradução prática:} ``Aceito 5\% de chance de estar errado.''

\medskip
\begin{tabular}{|c|c|c|l|}
\hline
\textbf{Confiança} & \textbf{Chance de erro ($\alpha$)} & \textbf{$Z$} & \textbf{Quando usar} \\
\hline
90\% & 10\% ($\alpha = 0{,}10$) & 1,645 & Pesquisas exploratórias, decisões reversíveis \\
95\% & 5\% ($\alpha = 0{,}05$) & 1,960 & Padrão em negócios e ciências sociais \\
99\% & 1\% ($\alpha = 0{,}01$) & 2,576 & Medicina, engenharia, decisões irreversíveis \\
\hline
\end{tabular}

\medskip
\textbf{Conexão fundamental:} o $\alpha = 0{,}05$ que define o nível de
confiança de 95\% é \textbf{o mesmo $\alpha = 0{,}05$} que vamos usar como
limiar de significância para o p-valor. São dois lados da mesma moeda.
\end{BoardBox}

Mas de onde vêm esses números $\alpha$ e $Z$? Não são arbitrários.
Eles vêm da \textbf{curva normal} (a famosa ``curva em sino'', ou
distribuição de Gauss) --- a mesma curva que vocês já viram quando
falamos de distribuição de dados.

\begin{SolvedBox}
\textbf{De onde vêm o $\alpha$ e o $Z$? --- A curva normal}

A lógica é visual.
Imagem a curva normal (sino) centrada na média da população:

\begin{enumerate}
  \item A curva mostra todas as médias possíveis que eu poderia obter
  se repetisse a pesquisa infinitas vezes.
  \item A \textbf{área total sob a curva} vale 1 (= 100\%).
  \item Quando eu digo ``95\% de confiança'', estou dizendo que quero
  que 95\% da área fique \textbf{dentro} do meu intervalo. Os 5\% que
  sobram (o $\alpha$) são divididos igualmente nas duas caudas:
  2{,}5\% à esquerda e 2{,}5\% à direita.
  \item O valor \textbf{$Z$} é o ponto na curva que separa a área
  central da área das caudas. Para 95\%, esse ponto é $Z = 1{,}96$:
  tudo entre $-1{,}96$ e $+1{,}96$ desvios-padrão da média corresponde
  a 95\% da área.
\end{enumerate}

\textbf{Na prática:}
\begin{itemize}
  \item \textbf{90\% de confiança} $\Rightarrow$ $\alpha = 0{,}10$
  $\Rightarrow$ sobram 5\% em cada cauda $\Rightarrow$ $Z = 1{,}645$.
  \item \textbf{95\% de confiança} $\Rightarrow$ $\alpha = 0{,}05$
  $\Rightarrow$ sobram 2{,}5\% em cada cauda $\Rightarrow$ $Z = 1{,}960$.
  \item \textbf{99\% de confiança} $\Rightarrow$ $\alpha = 0{,}01$
  $\Rightarrow$ sobram 0{,}5\% em cada cauda $\Rightarrow$ $Z = 2{,}576$.
\end{itemize}

Quanto mais confiança eu quero, mais largo precisa ser o intervalo
(maior $Z$), e portanto mais dados eu preciso coletar.

\textbf{Resumo em uma frase:}
$\alpha$ é a fatia de erro que eu aceito; $Z$ é a distância
em desvios-padrão que delimita essa fatia na curva normal.
Os dois são determinados \textbf{um pelo outro} --- escolheu
a confiança, $\alpha$ e $Z$ ficam definidos.
\end{SolvedBox}

Com o conceito de nível de confiança claro e sabendo de onde vêm $\alpha$
e $Z$, agora podemos apresentar a fórmula que transforma essas ideias
em um número concreto: o \textbf{tamanho mínimo da amostra}.
Na prática, é a resposta para ``quantas pessoas eu preciso entrevistar?''
ou ``quantas observações eu preciso coletar?''.

A fórmula calcula primeiro um valor chamado $n_0$.
O que é $n_0$? É o tamanho de amostra \textbf{ideal para uma população
infinita} --- ou seja, o número de observações que eu precisaria se a
população fosse tão grande que nunca ``acabasse''.
No mundo real a população tem um tamanho finito~$N$ (por exemplo,
2.000 clientes), então depois aplicamos uma \textbf{correção} que
reduz $n_0$ para o valor final $n$.
Pensem assim: $n_0$ é o ponto de partida (cenário mais exigente),
e a correção traz o número para a realidade.

\begin{FormulaBox}{Tamanho mínimo de amostra (para estimar uma média)}
\[
n_0 = \left(\frac{Z_{\alpha/2} \cdot \sigma}{e}\right)^2
\]
\begin{itemize}
  \item $n_0$ = tamanho de amostra para população infinita (ponto de partida)
  \item $Z_{\alpha/2}$ = valor $Z$ para o nível de confiança desejado (vem da curva normal)
  \item $\sigma$ = desvio-padrão estimado da população (de estudos anteriores ou piloto)
  \item $e$ = margem de erro tolerável (em unidades da variável)
\end{itemize}

\textbf{Para populações finitas} (quando $N$ é conhecido e pequeno):
\[
n = \frac{n_0}{1 + \dfrac{n_0 - 1}{N}}
\]
onde $N$ = tamanho total da população e $n$ = amostra final corrigida.

\textbf{Observação:} quando $N$ é muito grande ($N \gg n_0$), a correção
é desprezível e $n \approx n_0$ --- por isso $n_0$ sozinho já serve como
boa estimativa para populações grandes.
\end{FormulaBox}

A fórmula pode parecer abstrata à primeira vista, mas na prática é
bem direta.
Vamos aplicá-la a um caso concreto de uma rede varejista para ver
como cada peça se encaixa.

\begin{SolvedBox}
\textbf{Exemplo: quantos clientes preciso pesquisar?}

Uma rede varejista tem $N = 2.000$ clientes ativos e quer estimar o
ticket médio com 95\% de confiança e margem de erro de R\$~10.
Um estudo piloto indicou $\sigma \approx$ R\$~50.

\textbf{Passo 1 — Amostra sem correção:}
\[
n_0 = \left(\frac{1{,}96 \times 50}{10}\right)^2
    = \left(\frac{98}{10}\right)^2
    = 9{,}8^2
    = 96{,}04
    \approx 97
\]

\textbf{Passo 2 — Correção para população finita ($N = 2.000$):}
\[
n = \frac{97}{1 + \dfrac{97 - 1}{2000}}
  = \frac{97}{1 + 0{,}048}
  = \frac{97}{1{,}048}
  \approx 93
\]

\textbf{Resultado:} precisa de no mínimo \textbf{93 clientes}.

\textbf{E se quisesse 99\% de confiança?}
\[
n_0 = \left(\frac{2{,}576 \times 50}{10}\right)^2
    = 12{,}88^2
    \approx 166
\quad\Rightarrow\quad
n = \frac{166}{1 + \frac{165}{2000}} \approx 153
\]
Mais confiança $\Rightarrow$ mais dados necessários.
\end{SolvedBox}

Esse cálculo de amostra mínima vale para estimar médias, proporções
ou qualquer parâmetro.
Mas e no nosso contexto específico de \emph{correlação}?
A pergunta é ligeiramente diferente: \textit{``quantas observações preciso
para detectar uma correlação como significante?''}
A resposta depende do tamanho do próprio $r$ — correlações mais fracas
exigem mais dados para serem confirmadas.

\begin{BoardBox}
\textbf{E para correlação? Quantas observações eu preciso? (para o quadro)}

Para detectar se uma correlação é significante, o $n$ mínimo depende
do tamanho do efeito (o próprio $|r|$). Quanto menor a correlação que
você quer detectar, mais dados precisa.

\textbf{$|r|$ mínimo para significância a $\alpha = 0{,}05$ (bilateral):}

\medskip
\begin{tabular}{|r|c|l|}
\hline
\textbf{$n$} & \textbf{$|r|$ mínimo} & \textbf{Interpretação} \\
\hline
3  & 0,997 & Quase impossível detectar qualquer coisa \\
4  & 0,950 & Só correlação quase perfeita \\
5  & 0,878 & Precisa de correlação muito forte \\
6  & 0,811 & Ainda exige correlação forte \\
8  & 0,707 & Começa a ser viável para correlações fortes \\
10 & 0,632 & Mínimo razoável para análises exploratórias \\
15 & 0,514 & Detecta correlações moderadas a fortes \\
20 & 0,444 & Bom para investigações iniciais \\
30 & 0,361 & Detecta correlações moderadas \\
50 & 0,279 & Detecta correlações fracas a moderadas \\
100 & 0,197 & Detecta correlações fracas \\
500 & 0,088 & Detecta correlações muito fracas \\
\hline
\end{tabular}

\medskip
\textbf{Leitura da tabela:} com $n = 5$ e $r = 0{,}86$, o $r$ não
atinge o mínimo de 0,878 $\Rightarrow$ não é significante.
Com $n = 500$ e $r = 0{,}35$, o $r$ supera facilmente o
mínimo de 0,088 $\Rightarrow$ é significante.

\textbf{É por isso} que o cenário de 500 clientes com $r = 0{,}35$ é
mais confiável que 5 clientes com $r = 0{,}88$.
\end{BoardBox}

\begin{NoteBox}
\textbf{Regra prática para estudos com correlação:}
\begin{itemize}
  \item Para detectar correlação \textbf{forte} ($|r| \geq 0{,}7$): tenha pelo menos $n = 10$
  \item Para detectar correlação \textbf{moderada} ($|r| \geq 0{,}4$): tenha pelo menos $n = 25$
  \item Para detectar correlação \textbf{fraca} ($|r| \geq 0{,}2$): tenha pelo menos $n = 100$
  \item Na dúvida: \textbf{mais dados é sempre melhor}. Se possível, $n \geq 30$.
\end{itemize}
\end{NoteBox}

Agora que entendemos quantos dados precisamos e por que o tamanho da
amostra importa, vamos ao \textbf{teste formal} que responde à pergunta
central: ``essa correlação é estatisticamente significante?''.
A ferramenta para isso é o \textbf{p-valor} — um número entre 0 e 1
que mede a chance de o resultado que observamos ser puro acaso.
Vamos primeiro entender a lógica por trás, de forma intuitiva, antes
de entrar na fórmula.

\begin{BoardBox}
\textbf{O que é a Hipótese Nula ($H_0$)? (para o quadro)}

Antes de qualquer teste estatístico, precisamos definir o que estamos
testando. A ideia é a seguinte:

\begin{enumerate}
  \item Começamos \textbf{assumindo que não existe efeito nenhum} — essa
  suposição se chama \textbf{hipótese nula ($H_0$)}.
  \item Depois usamos os dados para tentar \textbf{derrubar} essa suposição.
  \item Se os dados forem fortes o suficiente para contradizer $H_0$, nós
  a \textbf{rejeitamos}. Se não forem, \textbf{não rejeitamos} — o que não
  significa que $H_0$ é verdade, apenas que não temos prova suficiente.
\end{enumerate}

\textbf{Analogia jurídica:} funciona como um julgamento.
O réu é considerado \textbf{inocente até que se prove o contrário}.
A ``inocência'' é a hipótese nula. As ``provas'' são os dados.
Se as provas forem fortes ($p < 0{,}05$), condenamos (rejeitamos $H_0$).
Se as provas forem fracas ($p \geq 0{,}05$), absolvemos — mas isso não
significa que o réu é de fato inocente, apenas que não houve evidência
suficiente para condenar.
\end{BoardBox}

Esse conceito de ``assumir que nada existe até provar o contrário'' é a
base de toda estatística inferencial.
Parece contraintuitivo no começo — por que começar negando? 
Porque é mais fácil provar que algo \emph{existe} do que provar que algo
\emph{não existe}.
Vejamos exemplos concretos para diferentes situações:

\begin{SolvedBox}
\textbf{Exemplos de $H_0$ e $H_1$ em diferentes contextos:}

\textbf{1. Correlação} (o nosso caso):
\begin{itemize}
  \item $H_0$: $\rho = 0$ — ``Não existe correlação entre X e Y na população.''
  \item $H_1$: $\rho \neq 0$ — ``Existe correlação entre X e Y.''
  \item \textit{Exemplo:} ``Horas de treino e performance não têm relação nenhuma''
  vs ``têm relação''.
\end{itemize}

\textbf{2. Comparação de médias:}
\begin{itemize}
  \item $H_0$: $\mu_A = \mu_B$ — ``A média dos dois grupos é igual.''
  \item $H_1$: $\mu_A \neq \mu_B$ — ``As médias são diferentes.''
  \item \textit{Exemplo:} ``O salário médio de homens e mulheres na empresa
  é o mesmo'' vs ``é diferente''.
\end{itemize}

\textbf{3. Proporção:}
\begin{itemize}
  \item $H_0$: $p = 0{,}50$ — ``A moeda é justa (50\% cara, 50\% coroa).''
  \item $H_1$: $p \neq 0{,}50$ — ``A moeda é viciada.''
  \item \textit{Exemplo:} ``A taxa de clique do botão vermelho e do azul é
  igual'' vs ``é diferente'' (teste A/B).
\end{itemize}

\textbf{Padrão:} $H_0$ é sempre ``nada acontece / não há diferença / não
há efeito''. $H_1$ é ``algo acontece''. O teste verifica se os dados
são fortes o suficiente para derrubar o $H_0$.
\end{SolvedBox}

Com a hipótese nula entendida, agora podemos falar do p-valor — que é
justamente a ferramenta que mede se os dados são fortes o suficiente
para rejeitar $H_0$.

\begin{SolvedBox}
\textbf{A lógica do p-valor — explicação intuitiva:}

O raciocínio completo é:
\begin{enumerate}
  \item \textbf{Definimos $H_0$:} ``A correlação na população é zero.
  Não existe associação real entre X e Y.''
  \item \textbf{Calculamos $r$} a partir da nossa amostra.
  \item \textbf{Perguntamos:} se $H_0$ fosse verdade (correlação zero na
  população), qual a probabilidade de observarmos um $r$ tão extremo
  quanto o que calculamos, \emph{por puro acaso}?
  \item \textbf{Essa probabilidade é o p-valor.}
\end{enumerate}

\textbf{Analogia da moeda:} imagine uma moeda que você suspeita ser viciada.
\begin{itemize}
  \item $H_0$: a moeda é honesta (50\% cara, 50\% coroa).
  \item Você joga 10 vezes e dá 7 caras.
  O p-valor é $\approx 0{,}17$ (17\% de chance de isso acontecer por acaso
  com moeda honesta). \textbf{Não rejeita $H_0$} — pode ser azar.
  \item Você joga 1.000 vezes e dá 700 caras.
  O p-valor é praticamente 0. \textbf{Rejeita $H_0$} — a moeda é viciada.
\end{itemize}

O mesmo vale para correlação: $r$ alto com poucos dados pode ser ``azar''.
$r$ alto com muitos dados quase certamente não é.

\textbf{Regra de decisão:}
\begin{itemize}
  \item \textbf{p $<$ 0,05:} ``É muito improvável ($<$ 5\%) que esse resultado
  seja acaso. Rejeitamos $H_0$ — a correlação é \textbf{estatisticamente
  significante}.''
  \item \textbf{p $\geq$ 0,05:} ``A chance de ser acaso é alta demais para
  concluir. Não rejeitamos $H_0$ — a correlação \textbf{não é significante}
  (o que não significa que não existe — apenas que não temos evidência
  suficiente).''
\end{itemize}

\textbf{Cuidado com a linguagem:} nunca dizemos ``aceitamos $H_0$''.
Dizemos ``não rejeitamos $H_0$'' — assim como no tribunal: absolver não
é o mesmo que provar inocência.
\end{SolvedBox}
A intuição ficou clara: o p-valor mede a probabilidade de o resultado
ser acaso.
Mas como calculá-lo na prática?
Para a correlação de Pearson, usamos um \textbf{teste~$t$} — uma fórmula
simples que transforma $r$ e $n$ em um valor que podemos comparar com
uma tabela de referência.

Mas antes de ver a fórmula, preciso que vocês entendam \emph{o que é}
esse valor $t$ e por que ele funciona como ``régua de surpresa''.

\begin{SolvedBox}
\textbf{O que o valor $t$ mede? — A régua de surpresa}

Pensem assim: o $t$ calculado mede \textbf{quão surpreendente} é o
resultado que obtivemos, \emph{assumindo que $H_0$ é verdade}
(ou seja, assumindo que a correlação real é zero).

\begin{itemize}
  \item $t$ \textbf{perto de 0:} ``Esse resultado não tem nada de especial.
  É exatamente o que esperaríamos por acaso.'' $\Rightarrow$ Não rejeita $H_0$.
  \item $t$ \textbf{moderado} (por exemplo 1,5): ``Hmm, ligeiramente
  incomum, mas ainda pode ser coincidência.'' $\Rightarrow$ Não rejeita $H_0$.
  \item $t$ \textbf{grande} (por exemplo 3,0): ``Isso seria muito raro
  se não houvesse correlação nenhuma. Provavelmente existe algo real aqui.''
  $\Rightarrow$ Rejeita $H_0$.
\end{itemize}

\textbf{Analogia:} imagine que alguém diz que é bom jogador de basquete.
\begin{itemize}
  \item Faz 6 de 10 cestas livres $\Rightarrow$ ``ok, pode ser sorte'' ($t$ baixo).
  \item Faz 9 de 10 cestas livres $\Rightarrow$ ``difícil ser só sorte'' ($t$ alto).
  \item Faz 95 de 100 cestas livres $\Rightarrow$ ``certeza que é bom'' ($t$ muito alto).
\end{itemize}

O número de tentativas (os dados) importa tanto quanto o resultado.
É por isso que a fórmula do $t$ combina o $r$ (o resultado) com o $n$
(a quantidade de dados).
\end{SolvedBox}

\begin{FormulaBox}{Teste $t$ para significância da correlação}
\[
t = r \cdot \sqrt{\frac{n - 2}{1 - r^2}}
\]
\begin{itemize}
  \item $r$ = coeficiente de correlação de Pearson calculado
  \item $n$ = número de observações (pares de dados)
  \item Graus de liberdade: $\text{gl} = n - 2$
  \item Compare $|t_{\text{calc}}|$ com $t_{\text{crítico}}(\alpha = 0{,}05;\; \text{gl})$
  \item Se $|t_{\text{calc}}| > t_{\text{crítico}}$ $\Rightarrow$ $r$ é significante
  \item O p-valor é obtido da tabela $t$ ou do software
\end{itemize}
\end{FormulaBox}

Vamos entender cada peça dessa fórmula para que não fique decoreba:

\begin{SolvedBox}
\textbf{Lendo a fórmula do $t$ peça por peça:}

\[
t = \underbrace{r}_{\text{sinal}} \cdot
    \sqrt{\frac{\overbrace{n - 2}^{\text{dados livres}}}
               {\underbrace{1 - r^2}_{\text{sobra não explicada}}}}
\]

\begin{itemize}
  \item \textbf{O $r$ na frente:} quanto maior o $r$ (a correlação),
  maior o $t$. Faz sentido --- correlação forte gera mais ``surpresa''.
  \item \textbf{O $n - 2$ no numerador:} são os graus de liberdade
  (dados livres). Quanto mais dados, maior o $t$. Faz sentido ---
  mais dados dão mais confiança.
  \item \textbf{O $1 - r^2$ no denominador:} é a parte da variação que
  a correlação \emph{não} explica. Se $r = 0{,}88$, então $r^2 = 0{,}77$
  e $1 - r^2 = 0{,}23$ --- sobram 23\% sem explicação. Quanto menor
  essa sobra (correlação mais forte), menor o denominador, e maior o $t$.
\end{itemize}

\textbf{Em resumo:} o $t$ fica grande quando a correlação é forte ($r$
alto), a amostra é grande ($n$ alto) e há pouca sobra inexplicada
($1 - r^2$ baixo). É exatamente quando devemos ter mais confiança.
\end{SolvedBox}

\begin{NoteBox}
\textbf{O limiar $\alpha = 0{,}05$:} é o nível de significância padrão —
aceita 5\% de chance de concluir que a correlação existe quando na verdade
não existe (erro tipo I). Em áreas de alto risco (medicina, engenharia)
pode-se exigir $\alpha = 0{,}01$ (apenas 1\% de chance de erro).
\end{NoteBox}

Agora, para usar a fórmula do teste $t$, precisamos consultar uma
\textbf{tabela} para saber qual é o ``limite'' ($t_{\text{crítico}}$) que
separa ``pode ser acaso'' de ``provavelmente não é acaso''.
Essa tabela depende de um conceito chamado \textbf{graus de liberdade}.
Soa complicado, mas é bem simples — vamos por partes.

\begin{BoardBox}
\textbf{O que são Graus de Liberdade ($\text{gl}$)? (para o quadro)}

Graus de liberdade é o número de valores que \textbf{podem variar
livremente} depois que fixamos alguma restrição (como a média).

\textbf{Analogia concreta:}
imagine que eu tenho 5 números e sei que a média deles é 10
(ou seja, a soma é 50). Quantos desses 5 números eu posso escolher
livremente?

\begin{itemize}
  \item Posso escolher o 1º: digamos 8. Faltam 42 para distribuir.
  \item Posso escolher o 2º: digamos 12. Faltam 30.
  \item Posso escolher o 3º: digamos 15. Faltam 15.
  \item Posso escolher o 4º: digamos 7. Faltam 8.
  \item O 5º \textbf{não tenho escolha}: tem que ser 8 (para a soma dar 50).
\end{itemize}

Ou seja: de 5 valores, apenas \textbf{4 são livres}. O último está
``preso'' pela restrição da média.
$\text{gl} = 5 - 1 = 4$.

\textbf{Para correlação:} temos $n$ pares de dados e estimamos 2 parâmetros
(as médias de X e de Y), então: $\text{gl} = n - 2$.

\medskip
\begin{tabular}{|c|c|l|}
\hline
\textbf{Situação} & \textbf{gl} & \textbf{Exemplo} \\
\hline
Correlação & $n - 2$ & 6 pares $\Rightarrow$ gl $= 4$ \\
Média de 1 grupo & $n - 1$ & 10 dados $\Rightarrow$ gl $= 9$ \\
Comparar 2 médias & $n_1 + n_2 - 2$ & 8+12 dados $\Rightarrow$ gl $= 18$ \\
\hline
\end{tabular}

\medskip
\textbf{Por que isso importa?} Porque com poucos graus de liberdade
(pouca informação livre nos dados), precisamos de um $t$ maior para
convencer que o resultado não é acaso. Com muitos graus de liberdade,
um $t$ menor já basta.
\end{BoardBox}

Percebam: quanto menos dados, menos graus de liberdade, e mais difícil
é provar que o resultado é significante. Faz sentido — com poucos dados,
precisamos de evidência mais forte.

Agora vamos à tabela em si. Mas antes, uma pergunta importante:
\textbf{de onde essa tabela veio?} Quem criou esses números e por quê?

\begin{SolvedBox}
\textbf{A história e a lógica por trás da tabela $t$ de Student}

Em 1908, William Sealy Gosset trabalhava como estatístico na cervejaria
Guinness, em Dublin. O problema dele era prático: a cervejaria fazia
testes de qualidade com \textbf{amostras pequenas} (poucos barris por lote).
A curva normal ($Z$) funcionava bem para amostras grandes, mas com
amostras pequenas os resultados ``oscilavam demais'' --- a curva normal
subestimava a incerteza.

Gosset percebeu que com poucos dados, a curva precisa ser \textbf{mais
achatada e com caudas mais largas} (mais espalhada) para refletir a
incerteza maior. Ele derivou matematicamente uma nova família de curvas
--- uma para cada valor de graus de liberdade. Publicou sob o pseudônimo
``Student'' (a Guinness não permitia publicações com nome real).

\textbf{A ideia visual:}
\begin{itemize}
  \item Imagine a curva normal (sino) que já conhecemos.
  \item Agora imagine que com \textbf{poucos dados} ($\text{gl}$ baixo),
  a curva fica mais \textbf{baixa e larga} --- as caudas ficam ``gordas'',
  significando que valores extremos são mais prováveis por acaso.
  \item Com \textbf{muitos dados} ($\text{gl}$ alto), a curva vai se
  afinando e ficando cada vez mais parecida com a normal.
  \item Com $\text{gl} = \infty$, a curva $t$ \textbf{se torna} a
  curva normal --- são a mesma coisa.
\end{itemize}

\textbf{Por que as caudas importam?}
O $t_{\text{crítico}}$ é o ponto que corta as caudas.
Se as caudas são gordas (poucos dados), esse ponto fica mais longe do
centro $\Rightarrow$ o $t_{\text{crítico}}$ é \textbf{maior}
$\Rightarrow$ é mais difícil rejeitar $H_0$.
Se as caudas são finas (muitos dados), o ponto fica mais perto
$\Rightarrow$ o $t_{\text{crítico}}$ é \textbf{menor}
$\Rightarrow$ é mais fácil rejeitar $H_0$.

\textbf{É por isso que a tabela existe:} cada linha é uma curva diferente,
e cada valor na tabela é o ponto de corte daquela curva específica.
Gosset calculou esses pontos de corte para que não precisássemos refazer
a matemática toda vez --- basta consultar a tabela.
\end{SolvedBox}

Resumindo: a tabela $t$ de Student é um \textbf{catálogo de limites de
surpresa}. Para cada quantidade de dados (graus de liberdade) e para
cada nível de exigência ($\alpha$), ela nos diz: ``o seu $t$ calculado
precisa ser maior que \emph{este valor} para você poder dizer que o
resultado não é acaso.''

Agora vamos ver a tabela em si e aprender a consultá-la.

\begin{BoardBox}
\textbf{A Tabela $t$ de Student — como ler (para o quadro)}

A tabela funciona como uma consulta em duas dimensões:
\begin{itemize}
  \item \textbf{Linha:} graus de liberdade ($\text{gl} = n - 2$ para correlação)
  \item \textbf{Coluna:} nível de significância ($\alpha$)
\end{itemize}

Nós vamos usar o caso mais comum: $\alpha = 0{,}05$ bilateral.

\textbf{O que significa ``bilateral''?}
Bilateral quer dizer que estamos testando se a correlação é
\textbf{diferente de zero} --- pode ser positiva ou negativa.
Por isso dividimos o $\alpha$ em duas metades: 2{,}5\% na cauda
esquerda (correlação negativa extrema) e 2{,}5\% na cauda direita
(correlação positiva extrema). Se testássemos apenas um lado
(``é positiva?''), seria \textit{unilateral} --- mas para correlação,
bilateral é o padrão.

O valor na interseção linha $\times$ coluna é o $t_{\text{crítico}}$.

\medskip
\textbf{Tabela $t$ crítico (bilateral) — recorte para o quadro:}

\medskip
\begin{center}
\begin{tabular}{|c|c|c|c|}
\hline
\textbf{gl} & $\alpha = 0{,}10$ & $\alpha = 0{,}05$ & $\alpha = 0{,}01$ \\
\hline
1  & 6,314 & 12,706 & 63,657 \\
2  & 2,920 & 4,303  & 9,925 \\
3  & 2,353 & 3,182  & 5,841 \\
4  & 2,132 & 2,776  & 4,604 \\
5  & 2,015 & 2,571  & 4,032 \\
6  & 1,943 & 2,447  & 3,707 \\
8  & 1,860 & 2,306  & 3,355 \\
10 & 1,812 & 2,228  & 3,169 \\
15 & 1,753 & 2,131  & 2,947 \\
20 & 1,725 & 2,086  & 2,845 \\
30 & 1,697 & 2,042  & 2,750 \\
60 & 1,671 & 2,000  & 2,660 \\
120 & 1,658 & 1,980  & 2,617 \\
$\infty$ & 1,645 & 1,960 & 2,576 \\
\hline
\end{tabular}
\end{center}

\medskip
\textbf{Como usar:}
\begin{enumerate}
  \item Calcule gl: $\text{gl} = n - 2$.
  \item Escolha a coluna de $\alpha$ (geralmente 0,05).
  \item Encontre o $t_{\text{crítico}}$ na interseção.
  \item Se $|t_{\text{calc}}| > t_{\text{crítico}}$: significante.
  Se $|t_{\text{calc}}| \leq t_{\text{crítico}}$: não significante.
\end{enumerate}

\textbf{Exemplo rápido de consulta:}
tenho $n = 8$ pares de dados e calculei $t = 2{,}50$.
\begin{enumerate}
  \item $\text{gl} = 8 - 2 = 6$.
  \item Coluna $\alpha = 0{,}05$.
  \item Na linha $\text{gl} = 6$: $t_{\text{crítico}} = 2{,}447$.
  \item $2{,}50 > 2{,}447$ $\Rightarrow$ \textbf{significante!}
  (por pouco, mas passou.)
\end{enumerate}
Se eu usasse $\alpha = 0{,}01$ (mais exigente), o $t_{\text{crítico}}$
seria $3{,}707$ --- e $2{,}50 < 3{,}707$, ou seja, \textbf{não seria
significante} ao nível de 1\%. O mesmo resultado pode ser significante a
5\% mas não a 1\% --- depende de quanto risco de erro você aceita.
\end{BoardBox}

Percebam dois padrões importantes na tabela:

Primeiro: à medida que os graus de liberdade \textbf{aumentam} (mais dados),
o $t_{\text{crítico}}$ \textbf{diminui}. Isso significa que com mais dados,
um valor de $t$ menor já é suficiente para provar significância.
Faz sentido: mais dados = mais confiança = menos exigência.

Segundo: a última linha diz $\text{gl} = \infty$.
O que é isso?

\begin{SolvedBox}
\textbf{O que significa $\text{gl} = \infty$ na tabela?}

Quando o número de graus de liberdade é muito grande (muitos dados),
a \textbf{distribuição $t$} se torna idêntica à \textbf{distribuição
normal} ($Z$) que já vimos antes.

Ou seja: $t_{\text{crítico}}(\text{gl} = \infty)$ é o próprio $Z$!

Vejam a coincidência (que não é coincidência):
\begin{itemize}
  \item $t_{\text{crítico}}(\text{gl} = \infty,\; \alpha = 0{,}10) = 1{,}645 = Z_{90\%}$
  \item $t_{\text{crítico}}(\text{gl} = \infty,\; \alpha = 0{,}05) = 1{,}960 = Z_{95\%}$
  \item $t_{\text{crítico}}(\text{gl} = \infty,\; \alpha = 0{,}01) = 2{,}576 = Z_{99\%}$
\end{itemize}

\textbf{São exatamente os mesmos valores $Z$} da tabela de nível de
confiança que vimos antes!

\textbf{Tradução prática:}
\begin{itemize}
  \item Com poucos dados ($n \leq 30$): use a tabela $t$ (mais exigente).
  \item Com muitos dados ($n > 120$): pode usar $Z$ — dá quase no mesmo.
  \item Na dúvida: a tabela $t$ sempre funciona (o $Z$ é um caso particular).
\end{itemize}
\end{SolvedBox}

Para fixar a leitura da tabela, vamos praticar com vários cenários
diferentes. Quero que vocês consigam consultar a tabela com naturalidade,
assim como consultamos uma tabela de preços.

\begin{BoardBox}
\textbf{Exercícios de consulta na tabela $t$ (para o quadro)}

Para cada cenário abaixo, identifique: (a) os graus de liberdade,
(b) o $t_{\text{crítico}}$ para $\alpha = 0{,}05$, e
(c) se o resultado é significante.

\medskip
\begin{tabular}{|r|c|c|c|c|l|}
\hline
\textbf{Cenário} & $n$ & $r$ & $t_{\text{calc}}$ & $\text{gl}$ & \textbf{Significante?} \\
\hline
A & 4  & 0,92 & 3,58 & ? & ? \\
B & 7  & 0,75 & 2,45 & ? & ? \\
C & 12 & 0,60 & 2,40 & ? & ? \\
D & 22 & 0,45 & 2,27 & ? & ? \\
E & 32 & 0,38 & 2,24 & ? & ? \\
\hline
\end{tabular}

\textbf{Resolução:}
\begin{itemize}
  \item \textbf{A:} $\text{gl} = 4 - 2 = 2$;
  $t_{\text{crítico}} = 4{,}303$.
  $3{,}58 < 4{,}303$ $\Rightarrow$ \textbf{NÃO significante}.
  (Mesmo com $r = 0{,}92$! Com apenas 4 dados, o limiar é altíssimo.)

  \item \textbf{B:} $\text{gl} = 7 - 2 = 5$;
  $t_{\text{crítico}} = 2{,}571$.
  $2{,}45 < 2{,}571$ $\Rightarrow$ \textbf{NÃO significante}.
  (Por muito pouco! Faltou um dado a mais para virar significante.)

  \item \textbf{C:} $\text{gl} = 12 - 2 = 10$;
  $t_{\text{crítico}} = 2{,}228$.
  $2{,}40 > 2{,}228$ $\Rightarrow$ \textbf{Significante}.
  (Aqui já tem dados suficientes.)

  \item \textbf{D:} $\text{gl} = 22 - 2 = 20$;
  $t_{\text{crítico}} = 2{,}086$.
  $2{,}27 > 2{,}086$ $\Rightarrow$ \textbf{Significante}.
  (Notem que o limiar caiu — com mais dados, ficou mais fácil.)

  \item \textbf{E:} $\text{gl} = 32 - 2 = 30$;
  $t_{\text{crítico}} = 2{,}042$.
  $2{,}24 > 2{,}042$ $\Rightarrow$ \textbf{Significante}.
  (Com 30 gl o limiar já é baixo; $r = 0{,}38$ modesto mas real.)
\end{itemize}

\textbf{Lição da tabela:} percebam como o cenário~A tem $r = 0{,}92$ e
não é significante, enquanto o cenário~E tem $r = 0{,}38$ e é.
É o tamanho da amostra (e portanto os graus de liberdade) que faz
toda a diferença. Mais dados $\Rightarrow$ $t_{\text{crítico}}$ menor
$\Rightarrow$ mais fácil de provar significância.
\end{BoardBox}

Com a tabela dominada, vamos agora aplicar tudo isso em dois exemplos
completos, passo a passo.
O primeiro mostra um caso em que a correlação \emph{é} significante;
o segundo, um caso em que \emph{não é} — mesmo com $r$ alto.
A diferença entre os dois vai ficar óbvia: é o tamanho da amostra.

\begin{SolvedBox}
\textbf{Cálculo passo a passo — exemplo significante:}

Um gestor quer saber se horas de treinamento e performance estão
correlacionadas na empresa. Dados de \textbf{6 funcionários}: $r = 0{,}88$.

\textbf{Passo 1 — Formular hipóteses:}
\begin{itemize}
  \item $H_0$: $\rho = 0$ (não há correlação na população)
  \item $H_1$: $\rho \neq 0$ (há correlação na população)
\end{itemize}

\textbf{Passo 2 — Calcular $t$:}
\[
t = 0{,}88 \cdot \sqrt{\frac{6 - 2}{1 - 0{,}88^2}}
  = 0{,}88 \cdot \sqrt{\frac{4}{1 - 0{,}7744}}
  = 0{,}88 \cdot \sqrt{\frac{4}{0{,}2256}}
  = 0{,}88 \cdot \sqrt{17{,}73}
  = 0{,}88 \times 4{,}21
  = 3{,}70
\]

\textbf{Passo 3 — Consultar $t_{\text{crítico}}$:}\\
Graus de liberdade: $\text{gl} = 6 - 2 = 4$.\\
Consultando a tabela $t$ do quadro, na linha $\text{gl} = 4$, coluna
$\alpha = 0{,}05$:\\
$t_{\text{crítico}} = 2{,}776$

\textbf{Passo 4 — Decisão:}\\
$|t_{\text{calc}}| = 3{,}70 > t_{\text{crítico}} = 2{,}776$

$\Rightarrow$ \textbf{Rejeitamos $H_0$}.
A correlação $r = 0{,}88$ é \textbf{estatisticamente significante} ao nível
de 5\%. O p-valor correspondente é $\approx 0{,}005$ ($< 0{,}05$).

\textbf{Interpretação em linguagem de negócio:}
``A associação entre horas de treinamento e performance é forte
($r = 0{,}88$) e estatisticamente significante ($p = 0{,}005$) — é muito
improvável que esse resultado seja mero acaso.''
\end{SolvedBox}

Agora vejam o que acontece quando o $r$ é alto, mas a amostra é muito
pequena.
O número muda completamente a conclusão:

\begin{SolvedBox}
\textbf{Cálculo passo a passo — exemplo NÃO significante:}

Uma consultora analisou a relação entre orçamento de marketing e NPS
em \textbf{4 filiais}: $r = 0{,}86$.

\textbf{Passo 1 — Hipóteses:} $H_0$: $\rho = 0$; $H_1$: $\rho \neq 0$

\textbf{Passo 2 — Calcular $t$:}
\[
t = 0{,}86 \cdot \sqrt{\frac{4 - 2}{1 - 0{,}86^2}}
  = 0{,}86 \cdot \sqrt{\frac{2}{1 - 0{,}7396}}
  = 0{,}86 \cdot \sqrt{\frac{2}{0{,}2604}}
  = 0{,}86 \cdot \sqrt{7{,}68}
  = 0{,}86 \times 2{,}77
  = 2{,}38
\]

\textbf{Passo 3 — Consultar $t_{\text{crítico}}$:}\\
$\text{gl} = 4 - 2 = 2$;
$t_{\text{crítico}}(\alpha = 0{,}05;\; \text{gl} = 2) = 4{,}303$

\textbf{Passo 4 — Decisão:}\\
$|t_{\text{calc}}| = 2{,}38 < t_{\text{crítico}} = 4{,}303$

$\Rightarrow$ \textbf{Não rejeitamos $H_0$}.
Mesmo com $r = 0{,}86$ — um valor visualmente alto — a correlação
\textbf{não é significante} ao nível de 5\%. O p-valor é $\approx 0{,}14$
($> 0{,}05$).

\textbf{Por quê?} Porque com apenas 4 observações, uma correlação de 0,86
pode facilmente surgir por acaso. Precisamos de mais dados para ter
confiança.

\textbf{Interpretação:}
``Os dados sugerem uma tendência positiva entre marketing e NPS
($r = 0{,}86$), mas com apenas 4 filiais, não temos evidência estatística
suficiente para afirmar que a associação é real ($p = 0{,}14$).
Recomendação: ampliar a análise para mais filiais antes de tomar decisão.''
\end{SolvedBox}

Esses dois exemplos deixam claro que $r$ e $n$ trabalham juntos.
Um $r$ alto com amostra pequena pode não significar nada;
um $r$ modesto com amostra grande pode ser altamente confiável.
Para consolidar, vamos organizar as combinações possíveis em um
quadro-resumo que vocês podem consultar sempre que precisarem
interpretar uma correlação.

\begin{BoardBox}
\textbf{Quadro-resumo: como reportar $r$ com significância (para o quadro)}

\medskip
\begin{tabular}{|c|c|c|l|}
\hline
\textbf{$r$} & \textbf{$n$} & \textbf{$p$} & \textbf{Interpretação} \\
\hline
$0{,}88$ & 6 & $0{,}005$ & Forte e significante — confiável \\
$0{,}86$ & 4 & $0{,}14$ & Forte mas NÃO significante — amostra pequena \\
$0{,}25$ & 500 & $< 0{,}001$ & Fraca mas significante — real, porém irrelevante \\
$0{,}95$ & 3 & $0{,}20$ & Fortíssimo mas NÃO significante — apenas 3 pontos \\
\hline
\end{tabular}

\medskip
\textbf{Regra prática:}
\begin{itemize}
  \item \textbf{Sempre} reporte $r$ \emph{e} $p$ juntos — nunca um sem o outro.
  \item $r$ alto + $p$ baixo: ``forte e confiável'' — ação justificada.
  \item $r$ alto + $p$ alto: ``parece forte mas pode ser acaso'' — precisa de mais dados.
  \item $r$ baixo + $p$ baixo: ``fraco mas real'' — a correlação existe, mas o efeito
  prático é pequeno.
  \item $r$ baixo + $p$ alto: ``nada encontrado'' — sem evidência de relação.
\end{itemize}
\end{BoardBox}

Com a teoria completa, vamos testar se vocês conseguem aplicar tudo
isso a um caso novo.
A questão a seguir mistura cenários com amostras diferentes — prestem
atenção ao $n$ e ao $p$ de cada unidade, porque a resposta muda
completamente dependendo desses valores.

\begin{SolvedBox}
\textbf{Questão para o Quadro — Interpretando $r$ e $p$}

Uma rede de academias pesquisou a relação entre frequência semanal
de treinos e perda de peso (kg) em 3 meses. Resultados:

\begin{center}
\begin{tabular}{|l|c|c|c|l|}
\hline
\textbf{Grupo} & \textbf{$n$} & \textbf{$r$} & \textbf{$p$} & \textbf{Significante?} \\
\hline
Unidade Centro & 8 & $0{,}82$ & $0{,}013$ & ? \\
Unidade Boa Vista & 120 & $0{,}21$ & $0{,}022$ & ? \\
Unidade Piedade & 5 & $0{,}90$ & $0{,}083$ & ? \\
\hline
\end{tabular}
\end{center}

Para cada unidade: (a) a correlação é significante? (b) qual sua
recomendação gerencial?

\textbf{Resolução:}
\begin{itemize}
  \item \textbf{Centro (n=8, r=0,82, p=0,013):} Significante (p $<$ 0,05).
  Correlação forte e confiável. Recomendação: investir em programas de
  aumento de frequência — há evidência de associação com resultados.
  \item \textbf{Boa Vista (n=120, r=0,21, p=0,022):} Significante
  (p $<$ 0,05), mas correlação fraca. A associação existe (amostra grande
  confirma), porém o efeito prático é pequeno — frequência sozinha explica
  pouco da perda de peso. Recomendação: investigar outros fatores (dieta,
  tipo de treino).
  \item \textbf{Piedade (n=5, r=0,90, p=0,083):} \textbf{NÃO significante}
  (p $>$ 0,05). Mesmo com $r$ alto, 5 observações são insuficientes para
  descartar o acaso. Recomendação: coletar mais dados antes de concluir.
\end{itemize}

\textbf{Lição:} significância e magnitude são coisas diferentes.
Um $r$ pequeno pode ser significante (com $n$ grande). Um $r$ grande
pode não ser significante (com $n$ pequeno). Reporte \textbf{ambos}.
\end{SolvedBox}

Na prática profissional, vocês não vão calcular o $t$ na mão — o
software faz isso.
Mas é fundamental entender o que está por trás, para interpretar os
resultados corretamente e não aceitar cegamente o que o computador
retorna.
Aqui vai como obter tudo em uma única linha de Python:

\begin{NoteBox}
\textbf{Em Python — como obter $r$ e $p$ juntos:}

\texttt{from scipy.stats import pearsonr}\\
\texttt{r, p = pearsonr(df['x'], df['y'])}\\
\texttt{print(f'r = \{r:.2f\}, p = \{p:.4f\}')}

O \texttt{df.corr()} do Pandas retorna apenas $r$ — \textbf{sem o p-valor}.
Para decisões que dependam de significância, use sempre o \texttt{pearsonr}
do \texttt{scipy}.
\end{NoteBox}

% ---------------------------------------------------------------------------
\subsection{21h30–21h40: Correlação não é Causalidade}
% ---------------------------------------------------------------------------

\begin{SolvedBox}
\textbf{O argumento de causalidade:}

Dizer que X \emph{causa} Y exige muito mais do que $r$ alto. Exige:
\begin{enumerate}
  \item \textbf{Precedência temporal:} X ocorre antes de Y.
  \item \textbf{Co-variação:} quando X muda, Y muda.
  \item \textbf{Ausência de variável confundidora:} não há Z que cause tanto X quanto Y
  e crie a correlação artificialmente.
  \item (Em pesquisa científica: experimento controlado ou desenho quasi-experimental.)
\end{enumerate}
\end{SolvedBox}

\begin{BoardBox}
\textbf{Correlações espúrias clássicas (para o quadro — debate em sala)}
\begin{itemize}
  \item \textbf{Chocolate × Prêmio Nobel:} países que mais consomem chocolate per capita
  têm mais ganhadores do Nobel. $r \approx 0{,}79$. Variável confundidora: riqueza
  (países ricos comem mais chocolate E investem mais em ciência).
  \item \textbf{Afogamentos × Sorvete:} mais afogamentos em meses de alta venda de sorvete.
  Causa real: \emph{verão} (calor leva as pessoas tanto à praia quanto a comer sorvete).
  \item \textbf{Storks × Birth rate (Países Baixos):} cidades com mais cegonhas têm mais
  nascimentos. Causa real: cidades rurais têm mais cegonhas E mais nascimentos (menos
  planejamento familiar). As cegonhas não trazem bebês — a ruralidade explica os dois.
  \item \textbf{Exemplo local:} ``Lojas que vendem guarda-chuvas têm mais acidentes de
  trânsito no mesmo dia.'' Causa real: chuva.
\end{itemize}

\textbf{Pergunta para a turma:} \textit{``Se você descobrir que o número de funcionários
de TI de uma empresa correlaciona positivamente com a receita, você recomendaria
contratar mais funcionários de TI para aumentar a receita? O que mais poderia
estar acontecendo?''}
\end{BoardBox}

\begin{SolvedBox}
\textbf{Correlação de Spearman (menção complementar):}

Quando as variáveis são ordinais ou a distribuição é muito assimétrica,
usamos a correlação de \textbf{Spearman} ($\rho$), que calcula Pearson usando
os \emph{postos} (ranks) dos valores, em vez dos valores absolutos.

Exemplo: ranking de satisfação (1–5 estrelas) × ranking de retenção de clientes.
Pearson seria inadequado porque as estrelas não têm intervalos iguais (salto de
1 para 2 estrelas não é o mesmo que 4 para 5).

Na prática em Python:
\begin{itemize}
  \item Pearson: \texttt{df[['x','y']].corr(method='pearson')}
  \item Spearman: \texttt{df[['x','y']].corr(method='spearman')}
\end{itemize}
\end{SolvedBox}

\begin{SolvedBox}
\textbf{Questão para o Quadro — Interpretando Scatter e Pearson}

Um analista de RH quer verificar se avaliação de desempenho e
salário se movem juntos. Dados de 6 funcionários:

\begin{center}
\begin{tabular}{|c|r|r|}
\hline
\textbf{Func.} & \textbf{Avaliação (0–100)} & \textbf{Salário (R\$ mil)} \\
\hline
A & 40 & 3{,}2 \\
B & 55 & 4{,}0 \\
C & 70 & 5{,}5 \\
D & 85 & 7{,}0 \\
E & 90 & 6{,}8 \\
F & 30 & 8{,}0 \\
\hline
\end{tabular}
\end{center}

O software calculou $r = 0{,}52$.

(a) Esboce mentalmente o scatter. Há outlier visual?\\
(b) Interprete o valor de $r$.\\
(c) O gerente diz: ``Então basta aumentar o salário para melhorar o
desempenho.'' Critique essa afirmação.

\textbf{Resolução:}\\
(a) Funcionários A--E seguem tendência positiva, mas \textbf{F é outlier
visual} (avaliação baixa, salário alto). Provavelmente um funcionário
com tempo de casa mas baixo desempenho recente.\\
(b) $r = 0{,}52$: correlação \textbf{moderada}. Existe associação positiva,
mas não é forte. O outlier F pode estar reduzindo o $r$.\\
(c) Correlação $\neq$ causalidade. Possíveis confundidores: (1)~anos de
experience \textrightarrow\ promovidos ganham mais E têm skills melhores;
(2)~área de atuação \textrightarrow\ áreas técnicas pagam mais e têm
métricas de desempenho diferentes. Aumentar salário sem contexto não
garante melhoria de desempenho.
\end{SolvedBox}

% =============================================================================
\section{Parte III — Prática em Dupla (21h40–21h55)}
% =============================================================================

\begin{BoardBox}
\textbf{Tabela 1 — Horas de Treinamento × Performance de Vendedor}

\medskip
\begin{tabular}{|c|l|r|r|}
\hline
\textbf{ID} & \textbf{Vendedor} & \textbf{Horas trein./mês} & \textbf{Perf. (0–100)} \\
\hline
01 & Vendedor A & 2  & 45 \\
02 & Vendedor B & 5  & 60 \\
03 & Vendedor C & 8  & 72 \\
04 & Vendedor D & 3  & 50 \\
05 & Vendedor E & 10 & 85 \\
06 & Vendedor F & 1  & 38 \\
07 & Vendedor G & 7  & 68 \\
08 & Vendedor H & 4  & 55 \\
09 & Vendedor I & 9  & 80 \\
10 & Vendedor J & 6  & 64 \\
\hline
\end{tabular}

\medskip
\textbf{Tarefas:}
\begin{enumerate}
  \item Calcule a média de X e Y.
  \item Desenhe (esquematicamente) o scatter: há padrão? Qual direção?
  \item A relação parece positiva, negativa ou ausente?
  \item Calcule $r$ (pelo menos 3 linhas do somatório para demonstrar raciocínio).
  \item Há outlier visual? Algum ponto que destoa do padrão?
\end{enumerate}

\textbf{Respostas esperadas:}
\begin{itemize}
  \item $\bar{x} = 5{,}5$; $\bar{y} = 61{,}7$
  \item Padrão: relação positiva forte — quanto mais treinamento, maior a performance
  \item $r \approx 0{,}99$ — correlação muito forte (conjunto criado quase linearmente)
  \item Sem outlier visual — todos os pontos seguem o padrão
  \item \textbf{Cuidado com causalidade:} pode ser que vendedores com melhor performance
  sejam escolhidos para mais treinamentos — não o contrário!
\end{itemize}
\end{BoardBox}

\begin{BoardBox}
\textbf{Tabela 2 — Temperatura × Vendas de Bebida Quente}

\medskip
\begin{tabular}{|c|r|r|}
\hline
\textbf{Dia} & \textbf{Temp. (°C)} & \textbf{Vendas (unidades)} \\
\hline
01 & 32 & 12  \\
02 & 28 & 18  \\
03 & 15 & 55  \\
04 & 10 & 88  \\
05 & 35 & 8   \\
06 & 22 & 35  \\
07 & 8  & 95  \\
08 & 18 & 48  \\
09 & 25 & 28  \\
10 & 12 & 70  \\
\hline
\end{tabular}

\medskip
\textbf{Tarefas:}
\begin{enumerate}
  \item Calcule a média de X e Y.
  \item Esboce o scatter. Qual é a direção da relação?
  \item Calcule $r$. A correlação é positiva ou negativa? Forte ou fraca?
  \item Isso constitui evidência de que temperatura \emph{causa} mudança nas vendas?
  O que mais poderia estar influenciando?
\end{enumerate}

\textbf{Respostas esperadas:}
\begin{itemize}
  \item $\bar{x} = 20{,}5$; $\bar{y} = 45{,}7$
  \item Padrão: relação negativa clara — quando temperatura sobe, vendas caem
  \item $r \approx -0{,}97$ — correlação negativa muito forte
  \item Causalidade plausível (frio $\to$ desejo de quente) mas não provada.
  Confundidor: feriados de inverno (mais movimento no café/restaurante em geral).
  Melhor framing: \textit{``temperatura é um forte preditor de vendas de bebida
  quente neste dataset — use-a para planejamento de estoque.''} Sem afirmar causa.
\end{itemize}
\end{BoardBox}

% =============================================================================
\section{Parte IV — Atividade Individual Final (21h55–22h)}
% =============================================================================

\begin{SolvedBox}
\textbf{Instruções:}
\begin{enumerate}
  \item Calcule médias de X e Y.
  \item Descreva o padrão do scatter (direção, intensidade, outlier visual).
  \item Calcule $r$ (mostre pelo menos o numerador e os dois somátorios do denominador).
  \item Interprete em linguagem de negócio (sem usar a palavra ``correlação'').
  \item Conclua: é possível afirmar causalidade com esses dados? Justifique.
\end{enumerate}
\textbf{Tempo:} 15 minutos.
\end{SolvedBox}

\begin{BoardBox}
\textbf{Tabela 3 — Investimento em Marketing × Novos Clientes}

\medskip
\begin{tabular}{|c|r|r|}
\hline
\textbf{Mês} & \textbf{Invest. Marketing (R\$ mil)} & \textbf{Novos Clientes} \\
\hline
Jan & 10 & 42  \\
Fev & 8  & 35  \\
Mar & 15 & 58  \\
Abr & 12 & 50  \\
Mai & 5  & 22  \\
Jun & 18 & 70  \\
Jul & 20 & 75  \\
Ago & 7  & 30  \\
Set & 14 & 53  \\
Out & 16 & 62  \\
\hline
\end{tabular}
\end{BoardBox}

\begin{SolvedBox}
\textbf{Gabarito (professor):}
\begin{itemize}
  \item $\bar{x} = 12{,}5$; $\bar{y} = 49{,}7$
  \item Padrão visual: relação positiva forte; sem outlier aparente
  \item $r \approx 0{,}996$ — correlação positiva muito forte
  \item \textbf{Linguagem de negócio:} ``Meses com maior investimento em marketing
  consistentemente apresentam maior captação de novos clientes — a relação é
  quase linear.''
  \item \textbf{Causalidade:} plausível, mas não provada. Possíveis confundidores:
  sazonalidade (meses com mais investimento coincidem com alta temporada?),
  eventos externos (lançamentos de produto?). Para estabelecer causalidade:
  experimento A/B controlando o orçamento de marketing.
\end{itemize}
\end{SolvedBox}

% ---------------------------------------------------------------------------
\section{Entrega — Trabalho 3 (lembrete)}
% ---------------------------------------------------------------------------

\begin{SolvedBox}
Entrega: \textbf{25/03/2026 às 23h59} via AVA.\\
Conteúdo: diagnóstico de dataset (missing + outliers + distribuição) + 1 scatter com análise
de correlação + Big Idea (dado → significado → decisão ou pergunta de investigação).
\end{SolvedBox}

% =============================================================================
\section{Revisão Semanal — Questão tipo Prova I}
% =============================================================================

\begin{NoteBox}
Esta semana trabalhamos os temas da \textbf{Questão 4} da prova:
correlação de Pearson, scatter plot e o cuidado com causalidade.
Vamos resolver uma questão no formato da prova para fixar.
\end{NoteBox}

\begin{SolvedBox}
\textbf{Cenário (estilo Q4 da Prova I):}

Um órgão estadual cruzou dados do ENEM de 120 escolas públicas e calculou
a correlação de Pearson entre \textit{nota média de matemática} e
\textit{índice socioeconômico (ISE)} das famílias: $r = 0{,}78$.

Um jornal publicou: \textit{``Estudo prova que renda familiar é a causa
das notas altas no ENEM.''}.

\textbf{Questão:} Analise a afirmação do jornal considerando:
\begin{enumerate}
  \item O que $r = 0{,}78$ realmente significa?
  \item Cite \textbf{dois confundidores} que poderiam explicar a correlação
  sem que renda cause diretamente a nota.
  \item Reescreva a manchete de forma técnica e responsável.
\end{enumerate}

\textbf{Resolução:}
\begin{enumerate}
  \item \textbf{Interpretação de $r$:} correlação positiva forte — escolas
  com ISE mais alto \textit{tendem} a ter notas mais altas. Isso indica
  \emph{associação}, não causalidade.
  
  \item \textbf{Confundidores plausíveis:}
  \begin{itemize}
    \item \textbf{Infraestrutura escolar:} famílias com renda mais alta
    vivem em regiões com escolas mais equipadas (labs, bibliotecas,
    professores mais qualificados). A infraestrutura impacta diretamente
    na aprendizagem.
    \item \textbf{Capital cultural:} famílias com maior ISE têm mais acesso
    a livros, viagens e estímulos intelectuais — isso influencia a nota
    independentemente da renda em si.
  \end{itemize}
  
  \item \textbf{Manchete corrigida:}
  \textit{``Dados do ENEM mostram forte associação entre índice
  socioeconômico e notas de matemática — especialistas alertam que
  múltiplos fatores podem explicar a relação.''}
\end{enumerate}

\textbf{Lição:} para afirmar causalidade, precisaríamos de (1) precedência
temporal, (2) co-variação e (3) controle de variáveis confundidoras
(experimento ou quasi-experimento). Uma correlação, por mais forte que
seja, é apenas o primeiro passo.
\end{SolvedBox}

% ---------------------------------------------------------------------------
\section{Materiais Preparados}
% ---------------------------------------------------------------------------
\begin{itemize}
  \item Slides: scatter plot animado; 4 padrões de correlação lado a lado
  \item Quarteto de Anscombe para debate em sala
  \item Tabelas 1, 2 e 3 impressas
  \item Site \texttt{tylervigen.com/spurious-correlations} para debate sobre causalidade
\end{itemize}

% ---------------------------------------------------------------------------
\section{Próximo passo}
% ---------------------------------------------------------------------------

Na \textbf{Semana 07} fechamos a análise bivariada com variáveis
\textbf{categóricas}: tabelas de contingência, proporções condicionais e
gráficos de barras agrupadas/empilhadas. Em seguida fazemos a revisão
completa para a \textbf{Prova I} (26 de março).

Perguntem aos alunos: \textit{``Como vocês descobririam se há correlação
entre a nota de satisfação do cliente (1–5 estrelas) e o prazo de entrega
(dias)? Seria Pearson ou Spearman? E como representariam visualmente?''}
